% Brief (tikz-diagram-craft).
% Reader question: what states can a claim enter, and what moves it between them?
% Claim: three states and five transitions, grounded in concrete developer API operations and SQLite semantics.
% Evidence: whitepaper/single-writer-kernel.tex sec:claim-lifecycle, alg:acquire.
\begin{figure}[H]
\centering
\pdfigurehue{pdcobalt}%
\begin{tikzpicture}[pd figure,x=1cm,y=1cm]
  \node[pd automaton state,minimum size=14.5mm] (free) at (1.3,0) {free};
  \node[pd automaton focus,minimum size=14.5mm] (held) at (5.7,0) {held};
  \node[pd automaton state,minimum size=14.5mm] (expired) at (10.1,0) {expired};
  \draw[pd initial] (-.3,0) -- (free.west);

  % acquire: the uniqueness decision that also tells the loser who holds it
  \draw[pd focus arrow] (free.east) -- (held.west)
    node[pd label,midway,above=2pt] {acquire};
  \node[pd focus badge] at (3.5,.95) {1};

  % re-claim: idempotent self-loop on held
  \draw[pd focus arrow] (held.115) .. controls +(115:1.1) and +(65:1.1) .. (held.65);
  \node[pd label,anchor=base] at (5.7,1.62) {re-claim};
  \node[pd focus badge] at (6.55,1.30) {2};

  % TTL elapses: held to expired, a plain rename, no badge
  \draw[pd arrow] (held.east) -- (expired.west)
    node[pd label,midway,above=2pt] {TTL elapses};

  % release: held back to free, routed below so it does not cross acquire
  \draw[pd arrow] (held.south) .. controls +(-100:1.0) and +(-80:1.0) .. (free.south);
  \node[pd label] at (3.5,-1.80) {release};
  \node[pd focus badge] at (3.5,-1.20) {4};

  % lazy sweep: expired back to free, on the next read, not on a schedule
  \draw[pd focus arrow,rounded corners=8pt]
    (expired.south) -- ++(0,-1.55) -- ($(free.south)+(0,-1.55)$) -- (free.south);
  \node[pd label,anchor=north] at (5.7,-2.40) {lazy sweep};
  \node[pd focus badge] at (5.7,-1.70) {3};

  \node[pd note,anchor=north west,text width=10.4cm] at (-.3,-2.85)
    {\textbf{No Queued State:} A losing acquire is denied immediately (\texttt{409 Conflict}); there is no in-memory buffer, wait queue, or fair scheduling. Losers must explicitly retry or back off.};
\end{tikzpicture}

\medskip
{\fontsize{4.8}{5.8}\selectfont
\begin{tabularx}{\linewidth}{@{}c >{\raggedright\arraybackslash}p{2.6cm} >{\raggedright\arraybackslash}p{3.1cm} >{\raggedright\arraybackslash}X@{}}
\toprule
\textbf{\#} & \textbf{API Trigger} & \textbf{Database Operation} & \textbf{Invariant \& Guarantee} \\
\midrule
1 & \texttt{POST /claims} \newline \texttt{\{"resource":"lib/auth.ts"\}} & \texttt{INSERT INTO claims ...} \newline On conflict: \texttt{SELECT holder} & Single atomic decision; winner gets \texttt{201 Created}; loser receives \texttt{409 Conflict} with current holder. \\
2 & \texttt{POST /claims/renew} \newline \texttt{\{"token":"tk\_94a2"\}} & \texttt{UPDATE claims SET expires\_at} \newline \texttt{WHERE holder = requester} & Fully idempotent; extends TTL without creating duplicate records or advancing the fencing epoch. \\
3 & Next read / acquire \newline \texttt{GET /claims} or \texttt{POST} & \texttt{DELETE FROM claims} \newline \texttt{WHERE expires\_at < now()} & Zero background threads; expired rows are swept lazily on demand before allocating new leases. \\
4 & \texttt{POST /claims/release} \newline \texttt{\{"token":"tk\_94a2"\}} & \texttt{DELETE FROM claims} \newline \texttt{WHERE holder = requester} & Immediate deterministic cleanup; excises the lock record upon task completion. \\
\bottomrule
\end{tabularx}}
\caption{The guarded claim lifecycle as a finite-state machine. Acquire is a single atomic insert against a SQLite uniqueness constraint; re-claim and release are idempotent; expiry moves a claim to \emph{expired}, from which the next read lazily sweeps it back to \emph{free}. There is no \emph{queued} state: losing callers receive \texttt{409 Conflict} and must manage their own backoff. [internal, alg:acquire]}
\label{fig:swk-claim-lifecycle}
\end{figure}
